\chapter{Coded Multi-Resource Systems with IDD}\label{ch:coded}

This chapter introduces channel coding and iterative detection and decoding (IDD) into the multi-resource PD-NOMA receiver. Motivated by the error-propagation problem observed in uncoded SIC receivers, the detector exchanges soft information with the FEC decoder and refines interference cancellation through iterations based on updated a priori information. This chapter develops MMSE-SIC-IDD and compares it with MMSE-PIC-IDD and MAP-IDD to clarify the performance-complexity tradeoff in coded multi-resource uplink NOMA.

\section{Brief introduction of IDD structure}

For coded systems, the detector is typically embedded in an iterative detection and decoding (IDD) receiver. In this structure, soft information is exchanged between the detector and the FEC decoder \cite{li2024turbo,sahin2024rsma,lin2024slotted}. The detector produces bit-level LLRs for the coded bits and passes them to the FEC decoder, while the decoder uses the code constraints to refine the reliability information and returns updated soft information to the detector. This detector-decoder exchange follows the IDD structure shown in Fig.~\ref{fig:idd-general}.

\figplaceholder{figures/idd_general_block.pdf}{General IDD receiver structure. The detector and FEC decoder exchange soft information through iterations based on a priori updates.}{fig:idd-general}

The iteration based on a priori update, denoted by $t$, is performed after all users have been decoded. The decoder outputs of all users are used to update their corresponding a priori information, and these updated a priori LLRs are then fed back to the detector. Through the a priori update iteration $t$, the updated a priori information is especially beneficial for both MMSE-SIC and MMSE-PIC. For MMSE-SIC, the updated information helps the stage-by-stage SIC process together with decoded-user feedback. For MMSE-PIC, it is even more critical because parallel cancellation relies directly on the a priori-based soft means of all interfering users.

\section{Coded system model}

Let each user transmit a coded bit block
\begin{equation}
    \bm{c}_i=[c_{i,1},c_{i,2},\ldots,c_{i,N_{c,i}}]^T .
\end{equation}
After modulation, the coded bits are grouped into $N_s$ modulation symbols. At symbol time $\ell$, the transmitted vector is
\begin{equation}
    \bm{x}_\ell=[x_{1,\ell},x_{2,\ell},\ldots,x_{U,\ell}]^T,
\end{equation}
and the received vector is
\begin{equation}
    \bm{y}_\ell=\bm{H}_\ell\bm{x}_\ell+\bm{n}_\ell.
    \label{eq:coded-symbol-model}
\end{equation}
The received block is $\bm{Y}=[\bm{y}_1,\bm{y}_2,\ldots,\bm{y}_{N_s}]$. In the block-fading simulations, $\bm{H}_1=\bm{H}_2=\cdots=\bm{H}_{N_s}$.

\section{MMSE-SIC-IDD}

\figplaceholder{figures/mmse_sic_idd_block.pdf}{MMSE-SIC-IDD receiver. At each SIC stage, a user is detected, decoded, and softly cancelled. After all users are processed, decoder extrinsic information is used for the next a priori update.}{fig:mmse-sic-idd-block}

\subsection{Receiver Procedure}

In this section, we introduce the procedure of the MMSE-SIC-IDD receiver. The key feature of the structure is that interference cancellation in the SIC process is performed after channel decoding rather than immediately after detection.

The received coded block is written as
\begin{equation}
    \bm{Y}
    =
    [\bm{y}_1,\bm{y}_2,\ldots,\bm{y}_{N_s}],
\end{equation}
where $N_s$ is the number of received coded symbols in the block. Let $\bm{H}_\ell$ denote the $\ell$th multiple-access channel and let $\bm{x}_\ell=[x_{1,\ell},x_{2,\ell},\ldots,x_{U,\ell}]^T$ denote the transmitted modulation-symbol vector at time index $\ell$. The $\ell$th column of $\bm{Y}$ is
\begin{equation}
    \bm{y}_\ell
    =
    \bm{H}_\ell\bm{x}_\ell+\bm{n}_\ell .
\end{equation}
Here $x_{i,\ell}$ denotes the transmitted modulation symbol of user $i$. For BPSK, $x_{i,\ell}=1-2c_{i,\ell}$; for higher-order modulation, $x_{i,\ell}$ is obtained by mapping a group of coded bits to a constellation point. The same IDD structure is used for all modulation orders.

Without loss of generality, we assume that the detection order has been determined according to the ordering criterion in \eqref{eq:post-sinr}, and that the users are re-numbered from $1$ to $U$ according to this order. After re-numbering,
\begin{equation}
    \Gamma^{\mathrm{post}}_1
    \geq
    \Gamma^{\mathrm{post}}_2
    \geq
    \cdots
    \geq
    \Gamma^{\mathrm{post}}_U ,
\end{equation}
and user $i$ is detected at the $i$th SIC stage. For user $i$, the MMSE-SIC detector first computes the detector extrinsic bit-level LLRs $\bm{L}_{\bm{c}_i}$ using \eqref{eq:idd-demapper-extrinsic-llr}. The detector also produces the coded soft-estimate vector
\begin{equation}
    \hat{\bm{s}}_i
    =
    [\hat{s}_{i,1},\hat{s}_{i,2},\ldots,\hat{s}_{i,N_s}]^T .
\end{equation}
The detector extrinsic LLRs are added to the a priori LLRs to form the decoder input LLRs $\bm{L}^{\mathrm{in}}_{\bm{c}_i}$. These decoder input LLRs are passed to the FEC decoder of user $i$. After decoding, the decoder produces the a posteriori LLRs $\bm{L}^{d}_{\bm{c}_i}$, which are used to generate the decoded soft symbol estimate $\hat{\bm{s}}^d_i$ according to the posterior-mean principle detailed in Section~\ref{sec:LLR calculation and soft estimate}. This decoded soft symbol estimate is then fed back to the MMSE-SIC detector for the next SIC stage, i.e., user $i+1$, where it is used to subtract the interference caused by user $i$. Intuitively, performing cancellation after decoding can improve the reliability of the cancelled signal and thus reduce error propagation in the SIC process.

Due to the successive decoding procedure, the extrinsic information of each user is collected asynchronously. Specifically, after decoding user $i$, the decoder extrinsic LLRs used for the next a priori update are obtained as
\begin{equation}
    \bm{L}^{e,d}_{\bm{c}_i}
    =
    \bm{L}^{d}_{\bm{c}_i}
    -
    \bm{L}^{\mathrm{in}}_{\bm{c}_i}
\end{equation}
for each user. A variance matrix is then used to collect the reliability information from all users. Once the extrinsic information of all users has been collected, it is used to update the a priori information for the next IDD iteration. In this way, the decoded extrinsic information of all users is provided as a priori information to the MMSE-SIC detector in the following IDD iteration.

\subsection{Soft Information and Variance}

Let $\mathcal{S}_i$ denote the modulation alphabet of user $i$, and let
$M_i=\log_2|\mathcal{S}_i|$ be the number of coded bits per modulation symbol. For a constellation point
$s\in\mathcal{S}_i$, its binary label is written as
\begin{equation}
    \bm{b}_i(s)
    =
    [b_{i,0}(s),b_{i,1}(s),\ldots,b_{i,M_i-1}(s)] .
\end{equation}
In this chapter, $x_{i,\ell}$ denotes the transmitted modulation symbol of user $i$ at time index $\ell$, while $s$ is used as a dummy constellation point when summing over $\mathcal{S}_i$.
Let $L^{(t)}_{c_i,\ell,m}$ denote the a priori LLR of the $m$th coded bit in the $\ell$th modulation symbol of user $i$ at a priori update iteration $t$, with
\begin{equation}
    \label{eq:a priori LLR}
    L^{(t)}_{c_i,\ell,m}
    =
    \log
    \frac{P^{a,(t)}(b_{i,\ell,m}=0)}
    {P^{a,(t)}(b_{i,\ell,m}=1)},
    \qquad
    L^{(0)}_{c_i,\ell,m}=0 .
\end{equation}
The corresponding bit probability is
\begin{equation}
    P^{a,(t)}(b_{i,\ell,m}=b)
    =
    \frac{1}
    {1+\exp\left[-(1-2b)L^{(t)}_{c_i,\ell,m}\right]},
    \qquad b\in\{0,1\}.
    \label{eq:bit-prior-from-llr}
\end{equation}
Assuming the interleaved coded bits are represented by independent bit priors at the demapper input, the a priori symbol probability is
\begin{equation}
    P^{a,(t)}_{i,\ell}(s)
    =
    \prod_{m=0}^{M_i-1}
    P^{a,(t)}(b_{i,\ell,m}=b_{i,m}(s)).
    \label{eq:symbol-prior-from-bit-priors}
\end{equation}
The soft symbol mean used for a priori update before detection is calculated according to the a priori symbol probability, that is,
\begin{equation}
    \mu^{(t)}_{i,\ell}
    =
    \E[x_{i,\ell}\mid \bm{L}^{(t)}_{c_i,\ell}]
    =
    \frac{
    \sum\limits_{s\in\mathcal{S}_i}
    s P^{a,(t)}_{i,\ell}(s)}
    {\sum\limits_{s\in\mathcal{S}_i}
    P^{a,(t)}_{i,\ell}(s)} ,
    \label{eq:idd-general-soft-mean}
\end{equation}
and the corresponding symbol variance is
\begin{equation}
    \sigma^{2~(t)}_{i,\ell}
    =
    \frac{
    \sum\limits_{s\in\mathcal{S}_i}
    \abs{s}^2 P^{a,(t)}_{i,\ell}(s)}
    {\sum\limits_{s\in\mathcal{S}_i}
    P^{a,(t)}_{i,\ell}(s)}
    -
    \abs{\mu^{(t)}_{i,\ell}}^2 .
    \label{eq:idd-general-soft-var}
\end{equation}
For BPSK with $0\mapsto +1$ and $1\mapsto -1$, \eqref{eq:idd-general-soft-mean} and \eqref{eq:idd-general-soft-var} reduce to
\begin{equation}
    \mu^{(t)}_{i,\ell}
    =
    \tanh\left(\frac{L^{(t)}_{c_i,\ell,0}}{2}\right),
    \qquad
    \sigma^{2~(t)}_{i,\ell}
    =
    1-\abs{\mu^{(t)}_{i,\ell}}^2 .
\end{equation}
The average variance of user $i$ over the coded block is
\begin{equation}
    \bar{v}^{(t)}_{i}
    =
    \frac{1}{N_s}\sum_{\ell=1}^{N_s}\sigma^{2~(t)}_{i,\ell},
\end{equation}
where $N_s$ is the number of modulation symbols in the coded block. The variance matrix is defined as
\begin{equation}
    \bm{V}
    =
    \operatorname{diag}(\bar{v}_1^{(t)},\bar{v}_2^{(t)},\ldots,\bar{v}_U^{(t)}),
    \qquad
    \bar{v}_i^{(t)}>0.
    \label{eq:variance-matrix}
\end{equation}
When decoder feedback is reliable, $\bar{v}_i^{(t)}$ becomes small and the receiver treats user $i$ as easier to cancel. When decoder feedback is uncertain, the variance remains large and the receiver avoids overconfident cancellation. This reliability weighting is essential in an IDD receiver.

\subsection{Variance-Aware MMSE Filtering}

When all users' extrinsic information from the FEC decoder has been collected, the block-averaged variance of each user is placed in the variance matrix $\bm{V}$ in \eqref{eq:variance-matrix}. The variance-aware MMSE filter used in the coded receiver is
\begin{equation}
    \label{eq:variance-aware-mmse}
    \tilde{\bm{w}}^H_k
    =
    \bar{v}_k^{(t)}\bm{h}_k^H
    \left(
    \bm{H}\bm{V}\bm{H}^H + 2\sigma^2\bm{I}
    \right)^{-1}.
\end{equation}
Compared with \eqref{eq:mmse-filter}, the covariance matrix includes $\bm{V}$, so users with reliable a priori information contribute less uncertainty to the MMSE filtering process.

\subsection{Decoder-Aided SIC Cancellation}

To simplify the expression in the following equations, we omit the subscript $\ell$, which denotes that the process is dealing with $\bm{y}_{\ell}$. The procedure then goes back to MMSE-SIC stage 1. User 1 benefits from the a priori information of users $2,\ldots,U$ by using
\begin{equation}
    \label{eq:idd-stage-one-residual}
    \bm{y}_{(1)}^{(t)}
    =
    \bm{y}
    -
    \sum_{j=2}^{U}
    \bm{h}_{j}\mu^{(t)}_{j}.
\end{equation}
At SIC stage $i>1$, the detection of user $i$ benefits from the FEC outputs of users $1,2,\ldots,i-1$ and from the a priori information of users $i+1,\ldots,U$. The residual vector for detecting user $i$ is constructed as
\begin{equation}
    \label{eq:idd-stage-i-residual}
    \bm{y}_{(i)}^{(t)}
    =
    \bm{y}
    -
    \sum_{j=1}^{i-1}
    \bm{h}_{j}\hat{s}^{d}_{j}
    -
    \sum_{j=i+1}^{U}
    \bm{h}_{j}\mu^{(t)}_{j}.
\end{equation}
The first summation cancels the users that have already been decoded in the current SIC pass by using the decoder-aided soft symbol estimates $\hat{s}^{d}_{j}$. The second summation cancels the users that have not yet been decoded in the current pass by using the a priori soft symbol means $\mu^{(t)}_{j}$. This is the coded extension of the uncoded subtraction in Chapter~\ref{ch:multi-resource}.

The cancellation term in \eqref{eq:idd-stage-i-residual} is different from the cancellation used in the uncoded MMSE-SIC detection. In the uncoded receiver, the cancelled signal is obtained directly from the detector output at the current detection stage. In the coded IDD receiver, the users $1,\ldots,i-1$ have already passed through the FEC decoder in the current SIC pass, so their decoder outputs provide more reliable soft estimates for cancellation. However, the cancellation is still not perfect. The residual uncertainty of the decoder-aided cancellation is represented by the cancellation residual variance $\bar{v}_c^{can}$. Its detailed calculation is given in Section~\ref{sec:LLR calculation and soft estimate}.

Similar to the uncoded MMSE-SIC procedure, the downsized channel matrix and the downsized variance matrix are used after each stage. The filter output for user $i$ is
\begin{equation}
    z^{(t)}_{i}
    =
    \tilde{\bm{w}}^H_i
    \bm{y}_{(i)}^{(t)},
\end{equation}
with equivalent scalar gain
\begin{equation}
    a^{(t)}_{i}
    =
    \tilde{\bm{w}}^H_i\bm{h}_i .
\end{equation}
The effective variance of the residual interference and noise is
\begin{equation}
    \sigma_{i,\mathrm{eff}}^{2~(t)}
    =
    \sum_{j>i}
    \bar{v}_j^{(t)}
    \abs{\tilde{\bm{w}}^H_i\bm{h}_j}^2
    +
    \sum_{c<i}
    \bar{v}_c^{can}
    \abs{\tilde{\bm{w}}^H_i\bm{h}_c}^2
    +
    2\sigma^2\norm{\tilde{\bm{w}}^H_i}^2 .
    \label{eq:idd-eff-var}
\end{equation}
The first term in \eqref{eq:idd-eff-var} comes from undecoded users whose interference is represented by the a priori variance $\bar{v}_j^{(t)}$. The second term is the cancellation residual caused by the already decoded users. The last term is the filtered noise variance. Therefore, unlike the uncoded MMSE-SIC detection, the coded receiver explicitly accounts for the residual uncertainty of decoder-aided cancellation.
Under the scalar approximation,
\begin{equation}
    z^{(t)}_{i}
    =
    a^{(t)}_{i}x_i
    +
    \eta^{(t)}_{i},
    \qquad
    \eta^{(t)}_{i}
    \sim
    \CN(0,\sigma_{i,\mathrm{eff}}^{2~(t)}).
\end{equation}

\subsection{LLR calculation and soft estimate}
 \label{sec:LLR calculation and soft estimate}
For the $q$th bit of user $i$, define
\begin{equation}
    \mathcal{S}_{i,q,b}
    =
    \{s\in\mathcal{S}_i:b_{i,q}(s)=b\},
    \qquad b\in\{0,1\}.
\end{equation}
After MMSE filtering and decoder-aided cancellation, the detector observes the scalar equivalent channel in \eqref{eq:idd-eff-var}. The detector LLR is therefore computed from $z_i^{(t)}$, the scalar gain $a_i^{(t)}$, the effective variance $\sigma_{i,\mathrm{eff}}^{2~(t)}$, and the a priori LLRs in \eqref{eq:a priori LLR}. The a priori symbol probability is given by \eqref{eq:symbol-prior-from-bit-priors}, where the bit probabilities are obtained from $L^{(t)}_{c_i,\ell,m}$ by the LLR-to-probability rule in \eqref{eq:bit-prior-from-llr}. The subscript $\ell$ is omitted in the following scalar detection equations for compactness.
The detector produces an extrinsic bit-level LLR by using the a priori information of the other bits in the same modulation symbol, but not the a priori information of the target bit itself:
\begin{equation}
    L_{c_{i,q}}
    =
    \log
    \frac{
    \sum\limits_{s\in\mathcal{S}_{i,q,0}}
    \exp\left(
    -\rho\frac{\abs{z^{(t)}_{i}-a^{(t)}_{i} s}^2}
    {\sigma_{i,\mathrm{eff}}^{2~(t)}}
    +
    \sum\limits_{\substack{m=0\\m\ne q}}^{M_i-1}
    \log P^{a,(t)}(b_{i,m}=b_{i,m}(s))
    \right)}
    {\sum\limits_{s\in\mathcal{S}_{i,q,1}}
    \exp\left(
    -\rho\frac{\abs{z^{(t)}_{i}-a^{(t)}_{i} s}^2}
    {\sigma_{i,\mathrm{eff}}^{2~(t)}}
    +
    \sum\limits_{\substack{m=0\\m\ne q}}^{M_i-1}
    \log P^{a,(t)}(b_{i,m}=b_{i,m}(s))
    \right)} .
    \label{eq:idd-demapper-extrinsic-llr}
\end{equation}
where $\rho$ is a detector LLR scaling factor. The ideal expression uses $\rho=1$, while the implementation uses $\rho=0.7$ to avoid overconfident detector output. The LLR actually passed to the FEC decoder is not only $L_{c_{i,q}}$, but the sum of the detector extrinsic LLR and the target-bit a priori LLR:
\begin{equation}
    L^{\mathrm{in}}_{c_i,\ell,q}
    =
    L_{c_{i,q}}
    +
    L^{(t)}_{c_i,\ell,q}.
    \label{eq:decoder-input-llr}
\end{equation}
For BPSK modulation, \eqref{eq:idd-demapper-extrinsic-llr} reduces to the detector extrinsic LLR
\begin{equation}
    L_{c_i}
    =
    \frac{4\rho\Re\{(a_i^{(t)})^*z_i^{(t)}\}}
    {\sigma_{i,\mathrm{eff}}^{2~(t)}}.
    \label{eq:idd-bpsk-llr}
\end{equation}
The corresponding BPSK decoder input is therefore $L^{\mathrm{in}}_{c_i,\ell}=L_{c_i}+L^{(t)}_{c_i,\ell,0}$.

The FEC decoder takes $\bm{L}^{\mathrm{in}}_{\bm{c}_i}$ as its input and produces posterior LLRs $\bm{L}^{d}_{\bm{c}_i}$. These posterior LLRs are used to form the decoder-aided soft estimate for SIC cancellation. The calculation follows the same bit-probability, symbol-probability, and posterior-mean rules in \eqref{eq:bit-prior-from-llr}, \eqref{eq:symbol-prior-from-bit-priors}, and \eqref{eq:idd-general-soft-mean}, with the a priori LLRs replaced by the decoder posterior LLRs. Therefore, for each symbol position $\ell$, the decoded soft estimate $\hat{s}^{d}_{i,\ell}$ is obtained from $\bm{L}^{d}_{\bm{c}_i}$ in the same way that $\mu^{(t)}_{i,\ell}$ is obtained from $\bm{L}^{(t)}_{c_i,\ell}$.

The only additional quantity required by decoder-aided SIC is the cancellation residual variance. Let $\sigma^{2,d}_{i,\ell}$ denote the symbol variance obtained by applying \eqref{eq:idd-general-soft-var} to the decoder posterior LLRs. Then the cancellation residual variance used in \eqref{eq:idd-eff-var} is the average posterior symbol uncertainty
\begin{equation}
    \bar{v}_i^{can}
    =
    \frac{1}{|\mathcal{N}_i|}
    \sum_{\ell\in\mathcal{N}_i}
    \sigma^{2,d}_{i,\ell},
    \label{eq:decoder-cancellation-residual-var}
\end{equation}
where $\mathcal{N}_i$ is the set of modulation-symbol positions occupied by user $i$. For BPSK modulation, the per-symbol decoded soft estimate and residual variance reduce to
\begin{equation}
    \hat{s}^{d}_{i,\ell}
    =
    \tanh\left(\frac{L^d_{c_i,\ell,0}}{2}\right),
    \qquad
    \sigma^{2,d}_{i,\ell}
    =
    1-\left(\hat{s}^{d}_{i,\ell}\right)^2 .
    \label{eq:bpsk-decoder-soft-estimate}
\end{equation}
The quantity $\bar{v}_i^{can}$ is the cancellation residual variance used in \eqref{eq:idd-eff-var}. It measures the remaining uncertainty after subtracting the decoder-aided soft estimate $\hat{s}^d_i$ from the received vector. A more reliable decoder output gives a smaller $\bar{v}_i^{can}$, which reduces the residual interference term in the following SIC stages.

After all users are processed, the detector a priori LLRs for the next a priori update are obtained from decoder extrinsic information:
\begin{equation}
    L^{(t+1)}_{c_i,\ell,q}
    =
    L^{d}_{c_{i,q}}
    -
    L^{\mathrm{in}}_{c_i,\ell,q}.
    \label{eq:decoder-extrinsic-feedback}
\end{equation}
Equivalently, in vector form, $\bm{L}^{e,d}_{\bm{c}_i}=\bm{L}^{d}_{\bm{c}_i}-\bm{L}^{\mathrm{in}}_{\bm{c}_i}$. The collection of all users' decoder extrinsic information is fed back to the detector as the a priori information for the next update. This a priori update is performed after all users have been decoded.

\subsection{Algorithm Summary}

\begin{algorithm}[t]
\caption{MMSE-SIC-IDD receiver for one coded block}
\label{alg:mmse-sic-idd}
\begin{algorithmic}[1]
\STATE Initialize the a priori LLRs $L^{(0)}_{c_i,\ell,m}=0$ for all users, coded-symbol positions, and bit positions.
\FOR{$t=0,1,\ldots,t_{\max}-1$}
    \STATE Compute $\mu^{(t)}_{i,\ell}$ and $\sigma^{2~(t)}_{i,\ell}$ for all users.
    \STATE Form the variance matrix $\bm{V}$ from the block-averaged variances.
    \STATE Determine the SIC order using post-SINR ordering and re-number the users as $1,\ldots,U$.
    \FOR{$i=1,2,\ldots,U$}
        \STATE Compute the variance-aware MMSE filter in \eqref{eq:variance-aware-mmse}.
        \STATE Build $\bm{y}^{(t)}_{(i)}$ using \eqref{eq:idd-stage-one-residual} or \eqref{eq:idd-stage-i-residual}.
        \STATE Generate detector extrinsic LLRs using \eqref{eq:idd-demapper-extrinsic-llr}.
        \STATE Form the decoder input LLRs using \eqref{eq:decoder-input-llr}.
        \STATE Decode user $i$ and obtain $\hat{s}^{d}_{i}$ and $\bar{v}_i^{can}$.
    \ENDFOR
    \STATE Set the next a priori LLRs using \eqref{eq:decoder-extrinsic-feedback}.
\ENDFOR
\STATE Make hard decisions from the final decoder outputs.
\end{algorithmic}
\end{algorithm}

\section{MMSE-PIC-IDD}

Parallel interference cancellation (PIC) cancels all other users in parallel instead of successively decoding one user at a time \cite{studer2011siso}. For user $k$, the residual vector is
\begin{equation}
    \bm{y}^{(t)}_{k,\ell}
    =
    \bm{y}_\ell
    -
    \sum_{j\ne k}\bm{h}_j\mu^{(t)}_{j,\ell}.
    \label{eq:pic-residual}
\end{equation}
The variance-aware MMSE filter follows the same form as \eqref{eq:variance-aware-mmse}, and the LLR calculation follows the MMSE-SIC procedure. PIC avoids SIC ordering and reduces decoding latency because users can be processed in parallel. However, PIC relies directly on the a priori-based soft means of all interfering users, so the quality of the a priori update is especially important. SIC, in contrast, performs ordered stage-by-stage cancellation and can additionally use decoded-user feedback from earlier stages in the same pass. This difference makes PIC and SIC complementary references: PIC represents a more parallel receiver whose cancellation quality is dominated by a priori reliability, whereas SIC represents a sequential receiver that can exploit current decoder outputs at the cost of ordering dependence and longer latency.

\section{MAP-IDD}

MAP-IDD provides a coded joint-detection baseline. In this receiver, the detector follows the maximum a posteriori (MAP) rule. The detector-decoder exchange follows the same IDD structure described previously: the detector generates bit-level LLRs, the FEC decoder refines them, and the decoder extrinsic information is fed back as the a priori information for the next update.

At symbol time $\ell$, the MAP detector jointly enumerates all multiuser symbol candidates. Let
$\bm{s}^{(q)}=[s^{(q)}_1,s^{(q)}_2,\ldots,s^{(q)}_U]^T$ denote the $q$th candidate vector, where $s^{(q)}_k\in\mathcal{S}_k$. The a priori LLRs $L^{(t)}_{c_k,\ell,m}$ are used in the detector by converting each candidate bit label into its log probability. Therefore, the MAP metric used for the $q$th candidate can be written as
\begin{equation}
    \Lambda_q
    =
    -\rho_{\mathrm{MAP}}
    \frac{\norm{\bm{y}_\ell-\bm{H}\bm{s}^{(q)}}^2}{2\sigma^2}
    +
    \sum_{k=1}^{U}
    \sum_{m=0}^{M_k-1}
    \log P^{a,(t)}
    \left(
    b_{k,\ell,m}
    =
    b_{k,m}\left(s^{(q)}_k\right)
    \right),
    \label{eq:map-idd-metric}
\end{equation}
where $\rho_{\mathrm{MAP}}$ is the detector metric scaling factor. In the implementation, $\rho_{\mathrm{MAP}}=0.7$, which corresponds to the scaling used before adding the a priori log-probability terms. Equivalently, the second term in \eqref{eq:map-idd-metric} is implemented by the same LLR-to-bit-probability rule as \eqref{eq:bit-prior-from-llr}.

For the $m$th bit of user $k$, define the candidate sets
\begin{equation}
    \mathcal{X}_{k,m,b}
    =
    \left\{
    \bm{s}^{(q)}:
    b_{k,m}\left(s^{(q)}_k\right)=b
    \right\},
    \qquad b\in\{0,1\}.
\end{equation}
The MAP detector output LLR is then
\begin{equation}
    L_{c_{k,m}}
    =
    \log
    \frac{
    \sum_{\bm{s}^{(q)}\in\mathcal{X}_{k,m,0}}
    \exp(\Lambda_q)}
    {
    \sum_{\bm{s}^{(q)}\in\mathcal{X}_{k,m,1}}
    \exp(\Lambda_q)} .
    \label{eq:map-idd-llr}
\end{equation}
This output is then passed to the FEC decoder, and the decoder feedback is used for the next a priori update as in the preceding IDD receiver descriptions. MAP-IDD provides strong BLER performance but retains exponential complexity in $U$ because it enumerates all multiuser symbol candidates. It is therefore useful as a benchmark for coded MMSE-SIC-IDD and MMSE-PIC-IDD.

\section{Coded Simulation Results}

This section evaluates the coded multi-resource NOMA system with IDD. The coded BLER performance is compared among MAP-IDD, MMSE-SIC-IDD with dynamic post-SINR ordering, and MMSE-PIC-IDD. Unless otherwise stated, the simulations use a $(1008,504)$ LDPC code over a block-fading channel, and the number of outer IDD iterations is denoted by $I_{\mathrm{out}}$. The purpose of the following results is to clarify how decoder feedback, SIC/PIC cancellation, system loading, and modulation order affect the coded receiver performance.
\begin{comment}
\begin{figure}[t]
    \centering
    \includegraphics[width=0.48\linewidth]{figures/coded_unequal_beta_sic_map.jpg}
    \hfill
    \includegraphics[width=0.48\linewidth]{figures/coded_unequal_beta_pic.jpg}
    \caption{Coded BLER performance for three users with unequal resource weights, where $R=2$ and $(\beta_1,\beta_2,\beta_3)=(1,0.1,0.01)$. The left figure compares MMSE-SIC-IDD with MAP-IDD, while the right figure shows MMSE-PIC-IDD with different outer iterations.}
    \label{fig:coded-unequal-beta}
\end{figure}
\end{comment}

According to Fig.~\ref{fig:coded-unequal-beta}, the BLER performance is strongly affected by the resource-weight imbalance. User 1, which has the largest resource weight, achieves the best performance, while user 3, which has the smallest resource weight, becomes much more difficult to decode. This result indicates that the effective reliability of each user is not only determined by the channel and code, but also by the resource-domain power distribution. For MMSE-SIC-IDD, one outer IDD iteration improves the BLER of all users and makes the performance very close to MAP-IDD. This shows that decoder-aided SIC cancellation can effectively reduce error propagation because the cancelled symbols are refined by the FEC decoder. In contrast, MMSE-PIC-IDD requires more IDD iterations to obtain reliable a priori information. When $I_{\mathrm{out}}$ is small, the parallel cancellation is inaccurate and the weak user remains interference-limited at high $E_b/N_0$.
\begin{comment}
\begin{figure}[t]
    \centering
    \includegraphics[width=0.82\linewidth]{figures/coded_user3_equal_beta.jpg}
    \caption{Coded BLER performance for the three-user case with equal resource weights and equal transmit powers.}
    \label{fig:coded-user3-equal}
\end{figure}
\end{comment}

Fig.~\ref{fig:coded-user3-equal} shows the BLER performance for the three-user case with equal resource weights. MAP-IDD gives the best performance because it performs joint detection over all possible superimposed signals. However, MMSE-SIC-IDD with one outer iteration approaches the MAP-IDD curve closely. The gain comes from the iterative exchange between the detector and the decoder: the decoder feedback improves the symbol reliability, and the variance-aware MMSE filter uses the updated a priori information to perform more reliable interference suppression. MMSE-PIC-IDD also improves as $I_{\mathrm{out}}$ increases, but the curve with $I_{\mathrm{out}}=0$ exhibits a clear error floor. Therefore, in the three-user case, MMSE-SIC-IDD provides a better performance-complexity tradeoff than MMSE-PIC-IDD when only a small number of IDD iterations is allowed.
\begin{comment}
\begin{figure}[t]
    \centering
    \includegraphics[width=0.82\linewidth]{figures/coded_user4_equal_beta.jpg}
    \caption{Coded BLER performance for the four-user case with equal resource weights and equal transmit powers.}
    \label{fig:coded-user4-equal}
\end{figure}
\end{comment}

Fig.~\ref{fig:coded-user4-equal} compares the receivers when the number of users is increased to four. Compared with Fig.~\ref{fig:coded-user3-equal}, the performance gap between MAP-IDD and the MMSE-based receivers becomes more visible because the multiuser interference is stronger. Nevertheless, MMSE-SIC-IDD still benefits from the IDD process, and the curve with $I_{\mathrm{out}}=1$ is clearly better than that with $I_{\mathrm{out}}=0$. It also outperforms MMSE-PIC-IDD in the simulated region. This result shows that, under moderate loading, the ordered and decoder-aided cancellation of MMSE-SIC-IDD is more effective than parallel cancellation based only on the current a priori estimates.
\begin{comment}
\begin{figure}[t]
    \centering
    \includegraphics[width=0.82\linewidth]{figures/coded_user6_equal_beta.jpg}
    \caption{Coded BLER performance for the six-user case with equal resource weights and equal transmit powers.}
    \label{fig:coded-user6-equal}
\end{figure}
\end{comment}

As shown in Fig.~\ref{fig:coded-user6-equal}, the performance gap becomes significant in the six-user case. MAP-IDD still maintains a steep BLER decrease as $E_b/N_0$ increases, whereas MMSE-SIC-IDD and MMSE-PIC-IDD both exhibit high-SNR error floors. This implies that, in a heavily loaded system, the dominant impairment is no longer thermal noise but residual multiuser interference after linear filtering and cancellation. Although the IDD iteration still improves MMSE-SIC-IDD, the improvement is not sufficient to close the gap to MAP-IDD. Therefore, MAP-IDD is more robust to the increase in the number of users, while the two MMSE-based receivers become more sensitive to system loading.
\begin{comment}
\begin{figure}[t]
    \centering
    \includegraphics[width=0.48\linewidth]{figures/coded_mixed_modulation_iter1.jpg}
    \hfill
    \includegraphics[width=0.48\linewidth]{figures/coded_mixed_modulation_iter0.jpg}
    \caption{User-wise coded BLER performance for mixed modulation with BPSK, QPSK, and 16QAM users. The results compare MMSE-SIC-IDD and MAP-IDD with different a priori update conditions.}
    \label{fig:coded-mixed-modulation}
\end{figure}
\end{comment}

Fig.~\ref{fig:coded-mixed-modulation} presents the coded BLER performance when different users employ different modulation orders. The BPSK user has the best BLER performance, followed by the QPSK user, while the 16QAM user is the most sensitive to noise and residual interference. This trend is expected because higher-order modulation has a smaller Euclidean distance between constellation points. In this mixed-modulation case, MAP-IDD still outperforms MMSE-SIC-IDD for all users. Moreover, comparing the results with and without a priori update shows that MAP-IDD can still benefit from decoder feedback even under block fading. The a priori information helps the MAP detector assign larger probabilities to symbol combinations that are more consistent with the decoder output, thereby improving the reliability of the final LLRs.

\section{Comparison and summary}

The simulation results show that the IDD structure is effective for coded multi-resource NOMA receivers, but the amount of gain depends on the detector type, system loading, resource-weight distribution, and modulation order. For MAP-IDD, the detector already provides highly reliable soft information by jointly evaluating the possible superimposed signals. Therefore, the additional gain from the a priori update is smaller than that of the MMSE-based receivers in simple BPSK cases. However, MAP-IDD still achieves the best BLER performance in all simulated scenarios and remains robust when the number of users or modulation order increases.

MMSE-SIC-IDD provides a favorable tradeoff between performance and complexity. By combining dynamic post-SINR ordering, variance-aware MMSE filtering, and decoder-aided cancellation, it can approach MAP-IDD in the three-user case and still performs well in the four-user case. The key advantage is that SIC cancellation is performed after channel decoding, so the cancelled symbols are more reliable than raw detector decisions. This reduces error propagation and improves the quality of the following SIC stages. However, MMSE-SIC-IDD is still based on linear filtering and sequential cancellation. When the system becomes heavily loaded, the accumulated residual interference causes a clear performance floor.

MMSE-PIC-IDD has the advantage of parallel processing and lower receiver latency because all users can be detected simultaneously. However, its cancellation quality strongly depends on the reliability of the a priori information. If the number of outer iterations is insufficient, the soft estimates used for parallel cancellation are inaccurate, and the receiver can suffer from an error floor. With more IDD iterations, MMSE-PIC-IDD improves, but it generally remains less reliable than MMSE-SIC-IDD in the moderate-loading cases considered here.

The unequal resource-weight results further show that weak users can become interference-limited if the resource-domain power distribution is highly imbalanced. The mixed-modulation results show that higher-order modulation users are more vulnerable to residual interference and require more reliable soft information. Overall, MAP-IDD is the most robust receiver but has the highest complexity, which grows rapidly with the number of users and modulation order. MMSE-SIC-IDD is a practical alternative for moderate system loading, while MMSE-PIC-IDD is more attractive when parallelism and latency are prioritized. These observations indicate that practical coded NOMA receiver design should jointly consider detector complexity, IDD iteration number, cancellation reliability, resource allocation, and modulation assignment.



\section{Receiver Tradeoffs}

Table~\ref{tab:receiver-tradeoff} summarizes the main receiver tradeoffs.

\begin{table}[t]
\centering
\caption{Qualitative comparison of uplink NOMA receiver families.}
\label{tab:receiver-tradeoff}
\begin{tabular}{p{0.24\linewidth}p{0.24\linewidth}p{0.24\linewidth}p{0.20\linewidth}}
\toprule
Feature & GDMA/HDS-CDMA & PD-NOMA soft SIC & Multi-resource PD-NOMA MMSE-SIC \\
\midrule
User separation & Complex channel gain and/or signature; magnitude and phase & Received-power difference; mainly magnitude & Resource-domain channel vectors; magnitude and phase through filtering \\
Detector type & Joint ML/MAP & Successive cancellation & Linear MMSE filtering plus SIC \\
Complexity & Exponential in number of users & Approximately linear in number of users & Approximately linear to polynomial, depending on matrix inversion size \\
Uncoded BER & Best among compared schemes & Sensitive to error propagation & Intermediate; improved by multiple resources \\
Coded BLER & Strong benchmark & Improved by decoder feedback if IDD is used & Competitive with IDD and variance-aware filtering \\
Latency & Higher due to joint enumeration & Low, but sequential & Low to medium; SIC is sequential, PIC can be parallel \\
System loading & Higher loading possible at higher complexity & Lower loading unless power disparity is strong & Medium loading with lower complexity than MAP \\
\bottomrule
\end{tabular}
\end{table}
